\documentclass[12pt]{article}
\providecommand{\CFIPaperDir}{}
% BEGIN Paper/preamble.tex
\usepackage[T1]{fontenc}
\usepackage{mathptmx}
\usepackage[margin=1in]{geometry}
\usepackage{setspace}
\usepackage[authoryear,round]{natbib}
\usepackage[hidelinks]{hyperref}
\urlstyle{same}
\doublespacing
\setlength{\emergencystretch}{1em}
\newcommand{\exhibitplaceholder}[3]{%
  \begin{quote}\small
    \textbf{#1 --- #2 (reserved).} #3
  \end{quote}%
}

% END Paper/preamble.tex
\hypersetup{pdftitle={Digital payments after remittance receipt},pdfauthor={}}
\title{Digital payments after remittance receipt:\\ Evidence from migrant women in Colombia}
\author{}
\date{Working draft --- not for submission}
\begin{document}
\maketitle
% BEGIN Paper/sections/abstract.tex
\begin{abstract}
Using a digital wallet in Colombia does not mean a woman can complete a transfer across borders or control how the money is used. We study migrant women's reported handling of remittances and participation in domestic digital payments, alongside existing interviews, focus groups, expert interviews and a five-country institutional comparison. Migrant accounts describe routine domestic use alongside reliance on intermediaries for some cross-border journeys, as well as variation by corridor. The study distinguishes reported use from the ability to open, operate and protect an account. It does not estimate the effect of Colombia's recent interoperable payments rollout.
\par\smallskip
\emph{Working-draft note: the quantitative findings sentence and final abstract await accepted exhibits.}
\end{abstract}
\noindent\textbf{Keywords:} digital payments; remittances; migrant women; financial inclusion; digital public infrastructure; Colombia.

% END Paper/sections/abstract.tex
% BEGIN Paper/sections/introduction.tex
\section{Introduction}
A woman can receive a remittance in a wallet and withdraw it all in cash. She can also receive cash and use digital payments for other transactions. A woman who runs a business using a local wallet may still need someone else to send money across a border. These are different forms of participation. Treating them as one measure of financial inclusion makes it hard to tell which part of the transaction works and which does not.

We study this distinction in Colombia. A survey of women with international migration experience separates the reported channel for the latest international remittance, the general disposition of remittance proceeds, and whether the respondent made or received a domestic digital payment or transfer in the previous 60 days. The latter two are the quantitative outcomes. They have different denominators and do not trace the same payment. Fee visibility, cash-out conditions, perceived safety, complaint outcomes and household decision arrangements provide context; they are not combined into a welfare measure.

Interviews and focus groups with Venezuelan migrant women document a contrast the survey outcomes cannot capture on their own. Women described using wallets and transfers for purchases, work and family payments within Colombia. For some transactions involving Venezuela, the local payment was instead the first step in a chain involving an intermediary, currency conversion and a separate confirmation from the recipient. Women also described different experiences in other corridors. The study therefore distinguishes digital ability, connection to a domestic payment service and the ability to complete an international transaction on terms the user can understand and verify.

This question complements work on earlier points of entry. \citet{barboni2025} study how perceived social acceptance shapes regularized Venezuelan migrants' willingness to engage with financial providers in Colombia. \citet{martinshields2022} distinguish access to technology from the use of digital government services in Bogotá. \citet{kabeer1999} separates resources, agency and achievements. We likewise distinguish access to a service from what people report doing with money and from who can decide how to use it. Neither a digital transaction nor a reported decision arrangement establishes a gain in agency.

Research on shared public digital systems distinguishes the functions of identity, payments and data exchange from their governance and adoption \citep{eavesrao2025}. Remittance studies likewise caution against reducing a transfer to a single rail: \citet{cirolia2022} describe how formal, informal and digital arrangements intersect in two Cape Town cases serving Congolese migrants. These studies do not tell us how common a specific practice is in Colombia. They help frame three questions for our evidence: what women say they do after receipt; which conditions accompany domestic use and retaining value digitally; and what can the migrant and institutional accounts tell us about where a transaction becomes difficult?

We first set the survey and fieldwork in Colombia's payments and documentation context, then separate the two engagement outcomes and their analytical bases. We read the migrant accounts before the expert diagnoses and regional cases. The discussion asks what a policy evaluation would need to measure beyond use.

% END Paper/sections/introduction.tex
% BEGIN Paper/sections/setting.tex
\section{Colombia's payments setting}
Bre-B entered full operation on 6 October 2025, following an implementation period in which private instant-payment initiatives and participating providers were already active \citep{banrep2026}. The central bank describes Bre-B as a service for domestic payments and transfers through participating financial institutions, not an application or an international remittance service \citep{banrepfaq}. Our survey interviews began after full operation, but the survey does not identify which rail processed a reported transaction. The calendar does not establish exposure or create a pre/post comparison.

Interoperability addresses whether providers can exchange a payment. It does not establish that a customer can enroll under her own name, navigate an interface without help, understand a charge or get a problem resolved. These are separate questions about implementation. Colombia's earlier PEP and later PPT were distinct regularization episodes; neither the recorded survey responses nor the authors' account of an onboarding difficulty verifies a participant's legal status or a provider's compliance.

The modern payment environment predates Bre-B. The framework for specialized electronic deposits and payments firms dates to 2014 \citep{law1735}. The 2021 protection statute creates a distinct temporary permit for eligible Venezuelan migrants \citep{decree216}. A legally recognized credential, a provider's technical ability to validate it, and a customer's continued use of her account are separate questions. The official history of Bre-B distinguishes earlier provider initiatives from full operation. Our recall windows apply to the survey questions, not to the launch date; they do not identify when an individual first used the rail.

% END Paper/sections/setting.tex
% BEGIN Paper/sections/data.tex
\section{Data and measurement}
\subsection{Survey}

The recorded survey interviews span 9 December 2025 to 17 February 2026. Fundación Capital recruited migrant women leaders, who reached potential participants through contacts and referrals and administered the survey. Recruitment was not restricted to former program participants. The intended population included adult women with international migration experience, including Colombian return migrants, with about three months' residence after arrival or return. The programmed screen used consent, gender, age and receipt of remittances by the respondent or someone in her household in the previous 12 months. The form did not test every intended recruitment condition. Household receipt does not establish that the respondent personally received or managed a transfer.

The revision retains the historical 423-record cohort and the original manual and coded exclusions. This is a nonprobability sample. Referral and household links were not systematically recorded, the cohort is not a count of verified distinct households, and there is no denominator for a response rate. Residence is recorded residence, not the interview site. The proposed cohort table (T1) describes the admitted records and the form-version partition; any additional geography panel has a separate source and disclosure check. A screening-flag diagnostic does not define a second set of exclusions.

The first outcome distinguishes respondents reporting any digital retention or use of remittance proceeds from those reporting withdrawal of all proceeds in cash. We also retain the five response categories in questionnaire order. Two choices describe amounts on the same side of the halfway point in different words. We do not reconstruct monetary shares or treat the choices as a strictly ranked scale. The latest receipt channel and the 12-month receipt frequency are separate context items; the latter preserves the combined response for nonapplicability or not remembering.

The second outcome is whether respondents report making or receiving a domestic digital payment or transfer in the previous 60 days. The denominator is valid yes/no answers, with refusal and other unscored states kept separate. Among affirmative domestic users, we retain the original five operation bins on the valid-bin base and report the response-state distribution on the affirmative-user base. Uncertainty is not a bin. Shared endpoints in the response labels cannot be resolved from the data. Fee visibility concerns the last operation in a separate 90-day window. The proposed T2 and F1 display the two outcomes without a common denominator or payment-flow arrows.

The working models use three survey measures. The access/basic-use measure combines reported phone access/use, internet-use frequency and running out of data or balance; some content is close to the outcomes. The second averages reported ability to send or receive money by phone and confidence in avoiding online fraud and suspicious messages. It measures neither tested skills nor confidence in making a transfer as a standalone construct. The third combines address documentation in the respondent's name and a phone line in her own name; these are two onboarding conditions, not account ownership or provider-verified eligibility. All specified components must be present. The optional three-item alternative adds app installation under its own applicability rule. Pair-complete and common-complete comparisons are distinct; Pearson correlations do not validate a scale.

\subsection{Interviews, focus groups and documentary comparison}

The original methods chapter reports 30 semistructured interviews with Venezuelan migrant women and six focus groups involving 47 participants, across Bogotá, Soacha, Cali and Medellín. These are reported activity counts, not a claim about the number of distinct people across activities. The qualitative analysis combined themes informed by the study's conceptual domains with themes developed from the interview and focus-group material. We retain the existing authored findings. We do not infer their frequency in the survey or assume participant-level linkage.

The original paper also draws on institutional and regional expert interviews and reviews regulatory, central-bank and policy documents. Experts describe implementation and institutional responsibilities; official sources document selected rules and system designs. Neither is evidence that a particular respondent could open an account, used a specific rail or resolved a complaint. The regional comparison examines counterpart systems, not comparable groups of women in five countries.

\exhibitplaceholder{T1}{Cohort and form versions}{Admitted total and version accounting. Screening states are a separate diagnostic, not exclusions. No geography enters from the cohort emitter alone.}

% END Paper/sections/data.tex
% BEGIN Paper/sections/quantitative.tex
\section{Quantitative approach and reporting}
We describe each outcome on its own valid-response base. Separate core logit models relate the two outcomes to the three measures, age, grouped education and grouped recorded residence. The reference categories are completed secondary education and Bogotá. The residual residence group is neither a fourth city nor a sampling cluster. Extensions add grouped occupation (employee or worker as reference) and original first-arrival calendar year, not reconstructed tenure.

The five model rows cover the two core models, their two extensions and a separate refit of the remittance core model on the shared outcome-observed core base. The design checks that the shared base equals the domestic core base before omitting a duplicate domestic fit. We distinguish intended eligibility from the estimator's sample. Completeness and valid inputs do not establish random missingness or justify treating model-specific exclusions as deletions from the cohort.

Each core contrast compares average fitted probabilities at the higher and lower specified values of a measure, with other covariates at their observed values and the averages over that model's estimation sample. The endpoints are 0.50 and 0.75 for access/basic use and for the transfer ability/online-safety measure, and 0.50 and 1.00 for the two onboarding conditions. These are different increments. Observed support at the endpoints does not establish joint support for all combinations used in standardization.

The proposed figure reports differences in percentage points with pointwise 95\% delta-method intervals from a working robust covariance that assumes independent records. The six raw and Holm-adjusted p-values concern coefficient-null tests in the two core models, not simultaneous confidence intervals for the plotted probability contrasts. The working covariance does not account for dependence through unrecorded referral or household links. Cross-sectional associations cannot distinguish skills acquired through use from skills that precede it, and adjustment does not remove unmeasured confounding. This is a reporting specification, not a preregistration claim.

\subsection{Survey results: exhibits reserved}
The observed outcome distributions come first, then the two core models and the separately specified checks. The slots below contain no estimates, counts from outputs, or rendered figures. See the technical supplement for field, denominator and source definitions.

\exhibitplaceholder{T2/F1}{Reported disposition and domestic digital participation}{Each panel uses its own valid-response base. Neither panel is a share of money or a stage in a transaction funnel.}
\exhibitplaceholder{QT1}{Working model status and sample sizes}{Show intended and estimator sample sizes and reasons for unavailability by model. E.5 hold.}
\exhibitplaceholder{QF1}{Adjusted probability differences}{Six fixed core contrasts in percentage points with pointwise 95\% intervals. E.5 hold.}
\exhibitplaceholder{QT2}{Six core coefficient tests}{Report raw and Holm-adjusted Wald-test probabilities, not simultaneous intervals. E.5 hold.}

% END Paper/sections/quantitative.tex
% BEGIN Paper/sections/migrant-findings.tex
\section{What migrant women describe}
\subsection{Domestic payments can work while a cross-border transfer still needs an intermediary}

Women described using wallets, bank transfers and QR payments to buy goods, receive customer payments and manage small businesses. In the accounts of sending money to Venezuela, a transfer from a Colombian wallet could start rather than finish the transaction. Women described asking trusted contacts or exchange operators for rates, transferring money locally, waiting for the other leg, and checking a receipt or asking relatives to confirm delivery. The legacy chapter also reports more direct or faster experiences in some other corridors; dependence on intermediaries was not uniform. These accounts concern different journeys and sometimes the opposite direction from the survey's inbound questions.

An intermediary could provide a service the local wallet did not: arrange the rate, coordinate delivery on the other side and confirm receipt. Some women used recommendations, repeat dealings and comparisons among contacts to decide whom to trust. These arrangements made transfers possible and sometimes offered useful reassurance. They also left women with work that did not end when the local transfer settled. Where charges and exchange-rate deductions were hard to separate, a fast local payment could still be part of a costly or time-consuming journey. These are reported experiences, not a price audit or a quantified loss of remittance value.

\subsection{Documents, practical help and control over the account}

Some women described trouble opening a product because they lacked a requested document; others said a document they held was not recognized by a provider or an app. Some relied on a relative's or partner's account. The same assistance could enable a transaction and limit privacy or direct control over it. Women described learning through use, comparing rates and helping others transact. The accounts tie these practices to livelihoods and to obligations such as paying for food, medicine or schooling across borders. Knowing how to move money and deciding where it should go do not guarantee control over the account or timely delivery.

\subsection{Taking precautions is not the same as having recourse}

The interviews and group discussions describe suspicious links, false confirmations and uncertainty about responsibility when a transfer was delayed or went wrong. Women reported checking requests and receipts, calling relatives and returning to known contacts. They could take these precautions and still find it hard to locate a complaint channel or know whether funds could be recovered. The legacy analysis distinguishes women facing barriers to formal entry, users who depend on intermediaries, digitally capable users constrained by available services, and women who transact more independently but still face corridor constraints. It treats security concerns as cross-cutting. These positions can overlap or change. They are not linked to survey respondents or assigned quantitative class labels.

% END Paper/sections/migrant-findings.tex
% BEGIN Paper/sections/expert-findings.tex
\section{What experts identify at the institutional level}
Experts distinguished formal recognition of a document from successful verification at a provider. They raised continuity when documentation changes, and practical disruptions such as limited data, unstable phone numbers and difficulties with app authentication. They proposed human assistance at onboarding and first use alongside simpler interfaces. This is an implementation diagnosis, not an estimate of the effect of training or support.

They also distinguished domestic payment interoperability from the cross-border work of verifying identity, converting currency, delivering funds at destination and handling disputes across organizations. A faster domestic payment can improve one leg without changing the others. In hybrid transactions the protection attached to that leg may not settle who is responsible for the rest of the transfer. Experts identified separate responsibilities for migration authorities, regulators, payment operators, providers and organizations offering assistance. The issue they raised was coordination across those responsibilities, not simply whether a rail or complaint channel existed. For a domestic Bre-B transaction, for example, the central bank directs customers to the participating financial institution's complaint channel \citep{banrepfaq}. That documented route does not say what occurred in any interviewee's case or establish a remedy for a cross-border dispute.

% END Paper/sections/expert-findings.tex
% BEGIN Paper/sections/benchmark.tex
\section{Institutional comparison: what can Colombia learn from its neighbours?}
The original benchmark uses the DPI Map \citep{ucldpimap}, operator and policy documents, expert interviews and Colombia's own findings. We keep the same five cases and the distinction between domestic system design and readiness to serve mobile or transnational users. The technical supplement records the evidence needed for the second question. A missing or unverified feature is not a low score.

Brazil's Pix illustrates a public operator's ability to set common rules and operate settlement and an alias directory \citep{bcbpix}. Peru offers a different example: the central bank phased interoperability in, starting by connecting Yape and Plin rather than asking users to move to a single national wallet \citep{bcrp2025}. Costa Rica's SINPE Móvil illustrates a simple phone-based transfer through participating institutions, but the number is linked to a local colón-denominated account \citep{bccrsinpe}. Mexico distinguishes SPEI interbank transfers from CoDi's mobile payment and collection interface \citep{banxicocodi,banxicospei}. Access to the underlying rail and use of a particular interface are different measures.

Colombia's Bre-B makes domestic interoperation available across participating providers. The comparisons help frame choices about operator responsibilities, phased connections, familiar interfaces and account requirements. They do not rank the five countries by migrant inclusion, verify acceptance of specific documents, or establish the effectiveness of cross-border protection. Any comparison of actual uptake, fees, recovery or international delivery needs its own population, period and source.

% END Paper/sections/benchmark.tex
% BEGIN Paper/sections/discussion.tex
\section{Discussion and policy implications}
The strands share a question, not a unit of observation. The survey distinguishes two reported uses. The interviews show how some women combine wallets, cash, trusted contacts and help from others. The expert findings describe why a product can work in one setting while another part of a transaction remains difficult. The benchmark documents what selected systems are designed to do. The strands complement each other without estimating a mechanism or validating one another's counts.

Three practical tests follow. First, can customers open and continue using an account under their own name when an accepted document changes, and get help when verification fails? Second, can a user see the total price and the amount due at destination before confirming a transfer? Third, if a transfer spans providers or an intermediary, can the user identify who handles the complaint, trace the payment and obtain a remedy? These tests separate usable access, transparent pricing and recourse from the number or speed of transactions. None presumes one universal barrier or treats cash use as failure.

The welfare implications concern things worth measuring: realized fees and exchange-rate spreads, travel or waiting time, delivery, losses and recovery, privacy, and control over spending. Some of these issues appear in the authored accounts; the survey's binary use outcomes cannot measure gains in them. Estimating welfare changes or ranking reforms would require outcomes and comparisons the present study does not provide. The contribution is to show why these questions should be asked at different stages of a payment.

% END Paper/sections/discussion.tex
% BEGIN Paper/sections/conclusion.tex
\section{Conclusion}
The migrant accounts make a distinction that is easy to miss when inclusion is measured by account access alone. A wallet can serve a household or small business in Colombia and still leave a woman dependent on someone else for a transfer abroad. Help from a relative or a trusted contact can make a transaction possible without giving the user control over the account or an effective remedy if the payment fails.

The practical question is what people can do at each step: enroll and maintain access, know what the recipient will receive, confirm delivery, and seek redress. Tracking those outcomes alongside use would help establish whether a change improves people's options or only makes one step faster. The present design does not estimate such gains.

\emph{Working-draft note: complete the quantitative conclusion and final declarations only after the joint evidentiary and release review.}

% END Paper/sections/conclusion.tex
% BEGIN Paper/references.tex
\begin{thebibliography}{99}
\bibitem[Banco Central de Costa Rica(n.d.)]{bccrsinpe}
Banco Central de Costa Rica. (n.d.). \href{https://www.bccr.fi.cr/cr/es/sistema-de-pagos/servicios-sinpe/servicios-dirigidos-a-personas.html}{\emph{Servicios dirigidos a personas: SINPE Móvil}}.

\bibitem[Banco Central de Reserva del Perú(2025)]{bcrp2025}
Banco Central de Reserva del Perú. (2025). \href{https://www.bcrp.gob.pe/docs/Sistema-Pagos/articulos/estrategia-de-interoperabilidad-2025-1.pdf}{\emph{Estrategia de interoperabilidad de los pagos minoristas en Perú: Principales avances al primer semestre del 2025}}.

\bibitem[Banco Central do Brasil(n.d.)]{bcbpix}
Banco Central do Brasil. (n.d.). \href{https://www.bcb.gov.br/en/financialstability/pixcentralbanksrole}{\emph{Central Bank's role in Pix}}.

\bibitem[Banco de la República(n.d.)]{banrepfaq}
Banco de la República. (n.d.). \href{https://www.banrep.gov.co/es/bre-b/preguntas-frecuentes}{\emph{Preguntas frecuentes de Bre-B}}.

\bibitem[Banco de la República(2026)]{banrep2026}
Banco de la República. (2026). \href{https://www.banrep.gov.co/es/blog/bre-b-breve-historia-sistema-interoperado-pagos-inmediatos}{\emph{Bre-B, una breve historia del sistema interoperado de pagos inmediatos}}.

\bibitem[Banco de México(n.d.-a)]{banxicocodi}
Banco de México. (n.d.-a). \href{https://www.banxico.org.mx/sistemas-de-pago/codi-cobro-digital-banco-me.html}{\emph{CoDi: Cobro Digital}}.

\bibitem[Banco de México(n.d.-b)]{banxicospei}
Banco de México. (n.d.-b). \href{https://www.banxico.org.mx/servicios/mi-spei_-transferencias-ban.html}{\emph{MI SPEI: Transferencias}}.

\bibitem[Barboni et~al.(2025)Barboni, de Roux, and Perez-Cardona]{barboni2025}
Barboni, G., de Roux, N., \& Perez-Cardona, S. (2025). \href{https://warwick.ac.uk/fac/soc/economics/research/centres/cage/manage/publications/wp745.2025.pdf}{\emph{Perceived social acceptance and migrants' financial inclusion}} (CAGE Working Paper No. 745/2025, February version). University of Warwick.

\bibitem[Cirolia et~al.(2022)Cirolia, Hall, and Nyamnjoh]{cirolia2022}
Cirolia, L. R., Hall, S., \& Nyamnjoh, H. (2022). \href{https://doi.org/10.1111/tran.12467}{Remittance micro-worlds and migrant infrastructure: Circulations, disruptions, and the movement of money}. \emph{Transactions of the Institute of British Geographers, 47}(1), 63--76.

\bibitem[Congreso de Colombia(2014)]{law1735}
Congreso de Colombia. (2014). \href{https://normograma.dian.gov.co/dian/compilacion/docs/ley_1735_2014.htm}{\emph{Ley 1735 de 2014}}.

\bibitem[Eaves and Rao(2025)]{eavesrao2025}
Eaves, D., \& Rao, K. (2025). \href{https://discovery.ucl.ac.uk/id/eprint/10210968/}{\emph{Digital public infrastructure: A framework for conceptualisation and measurement}} (Working Paper 2025-01). UCL Institute for Innovation and Public Purpose.

\bibitem[Kabeer(1999)]{kabeer1999}
Kabeer, N. (1999). \href{https://doi.org/10.1111/1467-7660.00125}{Resources, agency, achievements: Reflections on the measurement of women's empowerment}. \emph{Development and Change, 30}(3), 435--464.

\bibitem[Martin-Shields et~al.(2022)Martin-Shields, Camacho, Taborda, and Ruhe]{martinshields2022}
Martin-Shields, C. P., Camacho, S., Taborda, R., \& Ruhe, C. (2022). \href{https://doi.org/10.1002/poi3.280}{Digitalization and e-government in the lives of urban migrants: Evidence from Bogotá}. \emph{Policy \& Internet, 14}(2), 450--467.

\bibitem[Presidencia de la República de Colombia(2021)]{decree216}
Presidencia de la República de Colombia. (2021). \href{https://www.funcionpublica.gov.co/eva/gestornormativo/norma.php?i=159606}{\emph{Decreto 216 de 2021}}.

\bibitem[UCL Institute for Innovation and Public Purpose(n.d.)]{ucldpimap}
UCL Institute for Innovation and Public Purpose. (n.d.). \href{https://dpimap.org/methodology/}{\emph{Digital Public Infrastructure Map: Methodology}}.
\end{thebibliography}

% END Paper/references.tex
\end{document}
